% Generated by tools/gen_map.py (Russian original); English edition translated by hand.
\chapter{Mapping of Chapters to Logra Lessons}
\label{app:logramap}

The book is given in parentheses: ``Circuits'' and ``Verilog'' are books of the \emph{Learn} tab, ``Course'' is its guided lessons. The tables show which lesson
and which example (the \emph{Code} and \emph{Designer} tabs) correspond to the chapters of this book. The names were checked against the Logra
source code when this appendix was built.

\section*{Part I}
\begin{longtable}{@{}>{\raggedright\arraybackslash}p{0.24\textwidth}>{\raggedright\arraybackslash}p{0.38\textwidth}>{\raggedright\arraybackslash}p{0.26\textwidth}@{}}
\toprule Chapter & Logra lesson (book) & Example (\code{id}) \\ \midrule \endfirsthead
\toprule Chapter & Logra lesson (book) & Example (\code{id}) \\ \midrule \endhead
Signals & A bit: zero and one (Course); A logic gate (Course); A four-bit number (Course) & \code{button}, \code{not}, \code{blink} \\
Number Systems and Codes & Binary-coded decimal (Circuits); Gray code with a block (Circuits); The parity bit (Verilog) & \code{gray}, \code{parity}, \code{binary-bcd}, \code{hamming-encode} \\
Logic Gates and Boolean Algebra & NAND from blocks (Circuits); De Morgan's law (Circuits); XOR from four NANDs (Circuits) & \code{gates}, \code{inverse-gates}, \code{demorgan} \\
Combinational Circuits & Two out of three (Circuits); A multiplexer from NOT, AND, OR (Circuits); The ``two out of three'' vote (Verilog) & \code{logic}, \code{majority} \\
Minimizing Logic Functions & Equality: XNOR (Circuits); OR with three inputs (Circuits) & \code{implication}, \code{consensus} \\
Decoders and Encoders & Decoder (Circuits); Priority encoder (Circuits); 3-to-8 decoder (Verilog) & \code{decoder}, \code{priority} \\
Multiplexers and Demultiplexers & 2-to-1 multiplexer (Circuits); 4-to-1 multiplexer (Circuits); Demultiplexer (Circuits) & \code{mux}, \code{mux-four}, \code{demux-four} \\
Arithmetic Blocks & A half adder from gates (Circuits); A full adder (Circuits); Subtraction (Circuits); Less than: A < B (Circuits); Four-bit adder (Verilog) & \code{half-adder}, \code{full-adder}, \code{comparator} \\
Flip-Flops & One bit of memory (Course); SR flip-flop: set and reset (Circuits); D flip-flop: remember a bit (Circuits); Edge detector: one pulse per press (Circuits) & \code{edge-toggle} \\
Registers and Counters & A four-bit register (Circuits); A shift register as a block (Circuits); A counter as one block (Circuits); A counter with pause (Circuits) & \code{shift}, \code{counter}, \code{johnson} \\
Synchronous Circuits, FSMs and Memory & Synchronizer: tie the button to the clock (Circuits); A divide-by-2 frequency divider (Circuits); A ROM from an array (Verilog) & \code{async-toggle-sync}, \code{button-debounce}, \code{divide-by-three} \\
FSM Design and FIFO Queues & Build a traffic light (Course); The 101 detector (Circuits); A code lock (Circuits); A Moore output (Circuits) & \code{traffic}, \code{vending-fsm}, \code{sequence-1011}, \code{fifo-four} \\
Electrical Basics: CMOS, Outputs and Buses & An LED and current limiting (Circuits); A toggle switch (Circuits) & \code{button} \\
Digital and the Analog World: DAC, ADC, PWM & PWM: brightness by a number (Verilog) & \code{pwm}, \code{breathing-pwm}, \code{sigma-delta} \\
\bottomrule
\end{longtable}

\section*{Part II}
\begin{longtable}{@{}>{\raggedright\arraybackslash}p{0.24\textwidth}>{\raggedright\arraybackslash}p{0.38\textwidth}>{\raggedright\arraybackslash}p{0.26\textwidth}@{}}
\toprule Chapter & Logra lesson (book) & Example (\code{id}) \\ \midrule \endfirsthead
\toprule Chapter & Logra lesson (book) & Example (\code{id}) \\ \midrule \endhead
What Is an FPGA & What we have in hand (Course); 1. A tour of the workshop (Course) & \code{blink} \\
What Is Inside an FPGA & How long is one period? (Course) & \code{logic} \\
The Sipeed Tang Nano 20K Board & A button on a real board (Course); An LED that lights on zero (Verilog) & \code{button} \\
The Logra Environment & You are in a digital circuit workshop (Course); Where you have come to (Course) & \code{blink} \\
From Text to a Working Circuit & Hello! (Course); An LED and a button together (Verilog) & \code{button} \\
Verilog Basics & The first wire (Course); Concatenate bits (Course); Choice by condition (Course); The first memory cell (Course); A module inside a module (Course) & \code{gates}, \code{not} \\
From Theory to Verilog: Blocks of Part I & Full adder: accounting for the carry (Verilog); A counter with load (Verilog); Moore and Mealy: when the output reacts (Verilog) & \code{full-adder}, \code{decoder}, \code{traffic}, \code{fifo-four}, \code{hamming-syndrome} \\
Practice: Seven Projects & A running light (Verilog); Blinking: dividing the frequency by 8 (Verilog); Pressed: toggled (Verilog) & \code{gates}, \code{blink}, \code{chaser}, \code{button-debounce}, \code{traffic}, \code{breathing-pwm}, \code{mux} \\
Simulation & How long is one period? (Course) & \code{blink} \\
Clocking and Timing Analysis & A synchronizer of two flip-flops (Verilog); Reset immediately or on the clock (Verilog) & \code{async-toggle-sync}, \code{button-debounce} \\
UART, SPI, I2C & --- & \code{uart-tx}, \code{uart-rx}, \code{spi-loopback}, \code{i2c-pattern} \\
SDRAM Memory & An 8×4 RAM: write on the clock (Verilog); A ROM from an array (Verilog) & \code{ram-write-read}, \code{small-rom} \\
Building a Processor & A mini ALU for four bits (Verilog) & \code{alu}, \code{tiny-cpu} \\
Assignments for Independent Work & A button stopwatch (Verilog); Security: two sensors out of three (Verilog) & \code{stopwatch} \\
\bottomrule
\end{longtable}

\section*{The Running ``Security Panel'' Project}
\begin{longtable}{@{}>{\raggedright\arraybackslash}p{0.20\textwidth}>{\raggedright\arraybackslash}p{0.24\textwidth}>{\raggedright\arraybackslash}p{0.26\textwidth}>{\raggedright\arraybackslash}p{0.17\textwidth}@{}}
\toprule Stage & What we build & Lesson or example & Acceptance criterion \\ \midrule \endfirsthead
\toprule Stage & What we build & Lesson or example & Acceptance criterion \\ \midrule \endhead
Signals & A door sensor as a bit; an ``alarm'' LED on a button & A bit: zero and one (Course), \code{button} & The LED is lit while the button is pressed \\
Logic Gates & alarm = door AND (night OR guard) & AND with three inputs (Circuits), \code{gates} & The table on Waveforms matches the formula \\
Combinational Circuits and Minimization & Three sensors, a ``2 of 3'' vote, the minimal form & A ``2 of 3'' vote with a guard (Circuits); Security: two sensors out of three (Verilog), \code{majority} & All 8 combinations; at most one LUT \\
Decoders and Multiplexers & Choosing a zone and showing its number on seven segments & A digit on seven segments (Circuits), \code{decoder}, \code{mux} & The zone number is visible on the display \\
Arithmetic Blocks & A threshold: noise level A > B & Less than: A < B (Circuits), \code{comparator} & The alarm sounds only above the threshold \\
Codes & The lock code in BCD and a parity bit on transmission & Word parity check (Circuits), \code{parity} & An error in one bit is detected \\
Flip-Flops & Remembering the alarm, reset by a key & An alarm with memory (Circuits); An alarm with memory (Verilog), \code{edge-toggle} & The alarm holds after the signal goes away \\
Registers and Counters & An alarm-time counter, a divider for the beep & A pulse stretcher (Circuits), \code{counter} & The siren turns off after the set number of clocks \\
FSMs & OFF, ARMED, ALARM; entering the code 1-0-1 & A code lock (Circuits); A code lock: 1, 0, 1 (Verilog), \code{door-controller} & A wrong code does not disarm the panel \\
Electronics and PWM & A buzzer through PWM, pull-ups on the buttons, the board & A buzzer (Circuits); PWM: brightness by a number (Verilog), \code{pwm} & Works on the Tang Nano 20K board \\
\bottomrule
\end{longtable}
